\documentclass[preview]{standalone}

\usepackage{amsmath}
\usepackage{amssymb}
\usepackage{stellar}
\usepackage{definitions}
\usepackage{bettelini}

\begin{document}

\id{sequences-functions-functional-setting}
\genpage

\section{Practical criteria}

\begin{snippet}{uniform-convergence-supremum-formulation}
    A \sequence of \function[functions] \(f_n\colon S\fromto\realnumbers\)
    converges uniformly to \(f\colon S\fromto\realnumbers\) precisely when
    \[
        \lim_{n\to\infty}\sup_{x\in S}|f_n(x)-f(x)|=0.
    \]
    Thus the pointwise error is controlled by a single numerical sequence,
    independent of \(x\).
\end{snippet}

\begin{snippetproposition}{uniform-convergence-comparison-criterion}{Comparison criterion for uniform convergence}
    Let \(\{f_n\}\) and \(\{g_n\}\) be \sequence[sequences] of real-valued
    \function[functions] on \(S\). Suppose there exists \(C>0\) such that
    \[
        \sup_{x\in S}|f_n(x)-f_m(x)|
        \leq C\sup_{x\in S}|g_n(x)-g_m(x)|
    \]
    for every \(m,n\in\naturalnumbers\). If \(\{g_n\}\) converges uniformly
    on \(S\), then \(\{f_n\}\) also converges uniformly on \(S\).
\end{snippetproposition}

\begin{snippetproof}{uniform-convergence-comparison-criterion-proof}{uniform-convergence-comparison-criterion}{Comparison criterion for uniform convergence}
    Uniform convergence of \(\{g_n\}\) implies, by the uniform Cauchy
    criterion, that for every \(\varepsilon>0\) there exists
    \(N\in\naturalnumbers\) such that
    \[
        \sup_{x\in S}|g_n(x)-g_m(x)|<\frac{\varepsilon}{C}
    \]
    whenever \(m,n>N\). The assumed inequality then gives
    \[
        \sup_{x\in S}|f_n(x)-f_m(x)|<\varepsilon,
    \]
    so the uniform Cauchy criterion yields uniform convergence of
    \(\{f_n\}\).
\end{snippetproof}

\begin{snippetproposition}{uniform-convergence-moving-points-obstruction}{Moving-points obstruction to uniform convergence}
    Let \(f_n,f\colon S\fromto\realnumbers\). If there exists a
    \sequence \(\{x_n\}_{n\in\naturalnumbers}\subseteq S\) such that
    \[
        \limsup_{n\to\infty}|f_n(x_n)-f(x_n)|>0,
    \]
    then \(f_n\) does not converge uniformly to \(f\) on \(S\).
\end{snippetproposition}

\begin{snippetproof}{uniform-convergence-moving-points-obstruction-proof}{uniform-convergence-moving-points-obstruction}{Moving-points obstruction to uniform convergence}
    For every \(n\in\naturalnumbers\),
    \[
        \sup_{x\in S}|f_n(x)-f(x)|\geq|f_n(x_n)-f(x_n)|.
    \]
    The right-hand side does not tend to zero, so neither can the uniform
    error on the left.
\end{snippetproof}

\section{The space of bounded continuous functions}

\begin{snippetdefinition}{bounded-continuous-functions-space-definition}{Space of bounded continuous functions}
    Let \(S\subseteq\realnumbers\). The \emph{space of bounded continuous
    functions on \(S\)} is
    \[
        C_b(S)=\left\{
            f\colon S\fromto\realnumbers\suchthat
            f\text{ is continuous and bounded}
        \right\}.
    \]
\end{snippetdefinition}

\begin{snippet}{bounded-continuous-functions-uniform-norm}
    On \(C_b(S)\), define the uniform norm by
    \[
        \lVert f\rVert_\infty=\sup_{x\in S}|f(x)|.
    \]
    Its induced metric is
    \[
        d_\infty(f,g)=\lVert f-g\rVert_\infty.
    \]
    Convergence in this metric is exactly uniform convergence on \(S\).
\end{snippet}

\begin{snippettheorem}{bounded-continuous-functions-banach-theorem}{The space of bounded continuous functions is Banach}
    For every \(S\subseteq\realnumbers\), the space \(C_b(S)\) is a Banach
    space with respect to the uniform norm.
\end{snippettheorem}

\begin{snippetproof}{bounded-continuous-functions-banach-theorem-proof}{bounded-continuous-functions-banach-theorem}{The space of bounded continuous functions is Banach}
    Let \(\{f_n\}_{n\in\naturalnumbers}\) be a Cauchy \sequence in
    \(C_b(S)\), so
    \[
        \lVert f_n-f_m\rVert_\infty\longrightarrow0
        \qquad\text{as }m,n\to\infty.
    \]
    The uniform Cauchy criterion gives a \function
    \(f\colon S\fromto\realnumbers\) such that \(f_n\to f\) uniformly.
    A uniform limit of continuous functions is continuous, and a uniform
    limit of bounded functions is bounded. Hence \(f\in C_b(S)\), and
    \[
        \lVert f_n-f\rVert_\infty\longrightarrow0.
    \]
    Thus every Cauchy sequence converges in \(C_b(S)\).
\end{snippetproof}

\section{Polynomial approximation}

\begin{snippettheorem}{stone-weierstrass-real-compact-subset-theorem}{Stone--Weierstrass theorem on a compact subset of the real line}
    Let \(K\subseteq\realnumbers\) be compact. For every continuous
    \function \(f\colon K\fromto\realnumbers\), there exists a
    \sequence of real \polynomial[polynomials] \(\{P_n\}\) such that
    \(P_n\to f\) uniformly on \(K\). Equivalently, the restrictions to
    \(K\) of real polynomials are dense in \(C(K)=C_b(K)\) with respect to
    the uniform norm.
\end{snippettheorem}

\end{document}
